\section{The packet estimate and its phase class}\label{sec:packets}

We need a cubic bound for the samples obtained by synthesizing a
finite packet array and analyzing it at a larger scale. The input
array must remain arbitrary, since Section~\ref{sec:density} splits
it into coordinate subarrays. We first verify the precise phase
requirement of the cited packet theorem: uniform convexity alone
is not its stated hypothesis.

\subsection{Continued phases and local changes of coordinates}

The setup of \cite[Section~\PacketPhaseSection]{OpenAIPacket2026} starts with a
small spherical graph and extends its phase smoothly to
$\mathbb R^2$, with uniformly definite bounded Hessian and bounded
derivatives of every order at least two.  Its constants may depend on
this fixed continuation, the fixed bounded frequency region, and the
windows.  The continuation is made by cutting off the difference
from the quadratic Taylor polynomial.  We shall use continuations of
precisely this form: write
\[
 q_s(\xi)=-\sqrt{1-|\xi|^2},\qquad
 Q_s(\xi)=-1+\tfrac12|\xi|^2,
\]
and, for a sufficiently small fixed $\rho>0$, set
\begin{equation}\label{eq:spherical-continuation}
 \phi=Q_s+\zeta(q_s-Q_s),
 \qquad
 0\le\zeta\le1,\quad
 \zeta=1\ \text{on }B(0,\rho),\quad
 \operatorname{supp}\zeta\subset B(0,2\rho).
\end{equation}
Here $\zeta$ is smooth, the product is extended by zero outside
$B(0,2\rho)$, and the resulting phase must satisfy
Equation~\eqref{eq:ellipticity}.  No radiality is required of the cutoff.

\Needspace{5\baselineskip}
\begin{lemma}[Insertion of an elliptic graph]\label{lem:continuation}
Every germ of a smooth graph with positive definite Hessian can be
carried by an invertible affine transformation of $\mathbb R^3$ to a
graph agreeing, on an open neighborhood, with a phase of the form in
Equation~\eqref{eq:spherical-continuation}.  The phase, the neighborhood, and
the affine transformation can all be fixed independently of the
packet scales and of the exponents used below.
\end{lemma}

\begin{proof}
Choose a smooth cutoff $\zeta_0$ equal to one on $B(0,\rho)$,
supported in $B(0,2\rho)$, and identically $1/2$ on a neighborhood
$V$ of a point $\xi_1$ with $\rho<|\xi_1|<2\rho$.
It can be chosen by rescaling a fixed cutoff with an annular plateau.
The phase
\[
 q=Q_s+\zeta_0(q_s-Q_s)
\]
has Hessian $I+O(\rho^2)$: the remainder $q_s-Q_s$ and its
derivatives of orders zero, one, and two are respectively
$O(\rho^4)$, $O(\rho^3)$, and $O(\rho^2)$ on the support of the
cutoff.  Fix $\rho$ small enough that $q''$ is positive definite
everywhere.

After translating the target base point, let its height function be
$h_0$ near zero and put $H=h_0''(0)>0$ and $B=q''(\xi_1)>0$.
Choose an invertible real matrix $A$ with $A^{\mathsf T}HA=B$;
for example, $A=H^{-1/2}B^{1/2}$.  Define, near $\xi_1$,
\begin{align*}
 \psi(\xi)
 &=h_0\bigl(A(\xi-\xi_1)\bigr)+q(\xi_1)-h_0(0)\\
 &\quad+
 \bigl(\nabla q(\xi_1)-A^{\mathsf T}\nabla h_0(0)\bigr)
                \cdot(\xi-\xi_1).
\end{align*}
Thus $\psi$ and $q$ have the same value, gradient, and Hessian at
$\xi_1$.  This change of graph is induced by an invertible affine
map of ambient frequency space: its base part is
$x\mapsto \xi_1+A^{-1}x$, and its height part is addition of an
affine function of the new base coordinates.

Let $\vartheta\in C_c^\infty(B(0,2))$ equal one on $B(0,1)$ and
take $0\le\vartheta\le1$.  For a small fixed $\epsilon>0$ put
\[
 R_\epsilon(\xi)=
 \vartheta\bigl((\xi-\xi_1)/\epsilon\bigr)
                  \bigl(\psi(\xi)-q(\xi)\bigr).
\]
Choose $\epsilon$ so that its support is contained in $V$ and the
target graph is defined there.  Taylor's theorem and the agreement
of the two-jets give
\[
 \|D^jR_\epsilon\|_\infty\le C\epsilon^{3-j},
 \qquad j=0,1,2.
\]
Consequently $\phi=q+R_\epsilon$ remains uniformly convex for
sufficiently small $\epsilon$ and agrees with $\psi$ on
$B(\xi_1,\epsilon)$.  It equals $Q_s$ outside a compact set, so
all of its derivatives of order at least two are bounded.

To verify the required continuation form, observe that
$D=q_s-Q_s$ is strictly positive for $0<|\xi|<1$.
Shrink $V$, if necessary, so that $D\ge d_0>0$ there.  The function
$R_\epsilon/D$, extended by zero outside $V$, is smooth.  Set
\[
 \widetilde\zeta=\zeta_0+R_\epsilon/D.
\]
On the perturbation support, $\zeta_0=1/2$ and
$|R_\epsilon/D|\le C\epsilon^3/d_0$.  Choosing $\epsilon$
smaller makes this last quantity at most $1/4$.  Hence
$0\le\widetilde\zeta\le1$, with the same inner and outer support
properties as $\zeta_0$, and
$\phi=Q_s+\widetilde\zeta(q_s-Q_s)$.
Every choice was made using only the fixed target graph.
\end{proof}

In what follows, the phases to which the cited packet theorem is
applied are chosen as in Lemma~\ref{lem:continuation}.  Although the
inserted graph lies in the transition annulus, the packet theorem
allows arbitrary frequency labels in a fixed bounded region.  In
particular, those labels need not lie in the spherical region
$B(0,\rho)$: \cite[Definition~\PacketFamilyDefinition]{OpenAIPacket2026} imposes
no such restriction, and its single-step maps act on arbitrary
finite input arrays.  All windows meeting the inserted graph, and
the terminal windows used below, can therefore be included in one
fixed bounded frequency region.  The constants may depend on the
chosen phase, including its higher derivative bounds.

\subsection{Frames and the capacity norm}

We use the standard frequency and spatial localization of wave
packets; see Tao \cite[arXiv version~2, Section~4]{Tao2003} and the Fourier-series
construction of Guo, Oh, Wang, Wu, and Zhang
\cite[arXiv version~1, Section~4.1]{GOWWZ2025Packets}. Square windows give the exact
identities below, with the normalization needed for our sample sum.

Fix a dyadic $N\ge2$ and put $\delta=N^{-1}$ and
$U_\nu=-\nabla\phi(\nu)$.  For a dyadic factor $1\le m\le N$ let
\begin{equation}\label{eq:packet-scales}
 \theta=m\delta,\qquad r=m^{-2},\qquad
 w=r\theta=\delta/m.
\end{equation}
Partition $[0,1)$ into intervals $I$ of length $r$, with left
endpoints $t_I$.  At this scale, the indices are
\[
 (\nu,z,I),\qquad
 \nu\in\theta\mathbb Z^2,\qquad
 z\in D_{\mathrm f}^{-1}w\mathbb Z^2.
\]
Using the windows fixed in the local setup, define
\begin{align}
 (\mathsf P_{m,I}f)_{\nu,z}
 &=\int f(\xi)\chi_\theta^\nu(\xi)
              e\bigl(N^2(z\cdot\xi+t_I\phi(\xi))\bigr)\,d\xi,
              \label{eq:packet-analysis}\\
 \mathsf S_{m,I}d(\xi)
 &=(D_{\mathrm f}\theta)^{-2}
   \sum_{\nu,z}d_{\nu,z}\chi_\theta^\nu(\xi)
              e\bigl(-N^2(z\cdot\xi+t_I\phi(\xi))\bigr).
              \label{eq:packet-synthesis-general}
\end{align}
Synthesis is initially defined for finite arrays and then extended
continuously to $\ell^2$.

\begin{lemma}[Exact frame identities]\label{lem:packet-frame}
For every $f\in L^2(\mathbb R^2)$,
\begin{equation}\label{eq:packet-frame}
 \mathsf S_{m,I}\mathsf P_{m,I}f=f,
 \qquad
 \|\mathsf P_{m,I}f\|_{\ell^2}^2
       =(D_{\mathrm f}\theta)^2\|f\|_2^2,
 \qquad
 \|\mathsf S_{m,I}\|_{\ell^2\to L^2}
       =(D_{\mathrm f}\theta)^{-1}.
\end{equation}
\end{lemma}

\begin{proof}
Since $N^2w\theta=1$, the vectors $N^2z$ run through
$(D_{\mathrm f}\theta)^{-1}\mathbb Z^2$.  For each fixed
$\nu$, Equation~\eqref{eq:packet-analysis} is the Fourier
coefficient array of
$f\chi_\theta^\nu e(N^2t_I\phi)$ on a square of side
$D_{\mathrm f}\theta$.  Parseval and Fourier inversion give
\[
 \sum_z|(\mathsf P_{m,I}f)_{\nu,z}|^2
   =(D_{\mathrm f}\theta)^2\|f\chi_\theta^\nu\|_2^2
\]
and, after synthesis, $f(\chi_\theta^\nu)^2$.
Sum over $\nu$ and use $\sum_\nu(\chi_\theta^\nu)^2=1$.
The norm of synthesis follows from
$\mathsf S_{m,I}=(D_{\mathrm f}\theta)^{-2}\mathsf P_{m,I}^*$.
These identities first hold for finite partial Fourier sums in
$L^2$ and then for their limits.
\end{proof}

For $0<\kappa<1/10$, let $E$ be a centered closed ellipse with
ordered semiaxes $L\ge a>0$, of any orientation, and write
$W_\kappa(E)=L^{1+\kappa}a^{1-\kappa}$.
For a finite array $c$ in one bin at the scale
$(\theta,r,w)$ define
\begin{align}
 K_{\theta,r}(c)
 &=\sup_{\substack{E,u,v\\L\ge a\ge\theta}}
   \frac{\theta^2}{W_\kappa(E)}
   \sum_{\substack{U_\nu\in u+E\\z\in v+rE}}|c_{\nu,z}|^2,
                 \label{eq:packet-general-capacity}\\
 G_{\theta,r}(c)^3
 &=\left(w^2\sum_{\nu,z}|c_{\nu,z}|^2\right)
                          K_{\theta,r}(c)^{1/2}.
                 \label{eq:packet-gauge}
\end{align}
The translations $u,v\in\mathbb R^2$ are independent, and there
is no upper restriction on either semiaxis.  On each finite index
set let
\begin{align*}
 A_{\theta,r}(c)&=
 \inf\left\{\sum_{\ell=1}^qG_{\theta,r}(a_\ell):
       |c|\le\sum_{\ell=1}^q|a_\ell|,\ q<\infty\right\},\\
 \mathcal A_m(c)^3&=r\sum_I A_{\theta,r}(c_I)^3.
\end{align*}
The domination in this infimum is coordinatewise.

\Needspace{5\baselineskip}
\begin{lemma}[The atomic norm controls the sample sum]
\label{lem:packet-atomic}
The function $A_{\theta,r}$ is a norm on every finite coordinate
space, is monotone under absolute coordinatewise domination, and
is contracted by coordinate masks.  Moreover,
\begin{equation}\label{eq:packet-atomic-lower}
 rw^2\sum_{\nu,z,I}|c_{\nu,z,I}|^3\le\mathcal A_m(c)^3.
\end{equation}
\end{lemma}

\begin{proof}
The disk of radius $\theta$, with its two translations chosen
independently, contains any prescribed index in both tests.
Its weight is $\theta^2$, so
$K_{\theta,r}(c)\ge\|c\|_\infty^2$ and
\[
 G_{\theta,r}(c)^3
 \ge w^2\|c\|_2^2\|c\|_\infty
 \ge w^2\sum_{\nu,z}|c_{\nu,z}|^3.
\]
For every atomic domination of $c$, the triangle inequality in
$\ell^3$ now gives
\[
 w^{2/3}\|c\|_3\le\sum_\ell G_{\theta,r}(a_\ell).
\]
Taking the infimum and summing its cube over bins with weight $r$
proves Equation~\eqref{eq:packet-atomic-lower}.  A singleton has
cost $w^{2/3}|c_{\nu,z}|$, so the infimum is finite; the displayed
lower bound makes it positive away from zero.  Homogeneity,
subadditivity, and absolute monotonicity follow directly from the
definition, as does contraction under masks.
\end{proof}

At the initial scale $m=1$ there is one bin, $[0,1)$, and we use
the notation
\begin{equation}\label{eq:packet-capacity}
 \mathcal K_\kappa(d)=K_{\delta,1}(d),
 \qquad
 m(d)=\delta^2\sum_{\nu,z}|d_{\nu,z}|^2.
\end{equation}
In particular,
$G_{\delta,1}(d)^3=m(d)\mathcal K_\kappa(d)^{1/2}$.

\subsection{The cited theorem and its specialization}

A single-step map $T_m$ synthesizes at factor $1$, analyzes in
every factor-$m$ time bin, and then applies any coordinate mask.
It uses the full matrix between its input and retained output
indices.  The admissible complexity convention of
\cite[Definition~\PacketFamilyDefinition]{OpenAIPacket2026} fixes a finite $B$
and a bounded frequency-label region, requires
$m\le N\le m^B$ and position labels of size at most $m^B$, and
allows arbitrary finite arrays on those indices.  The phase and
window constants are fixed throughout this class.

\begin{theorem}[Packet propagation
{\cite[Theorem~\PacketTheoremNumber, Equation~(\PacketBoundEquation)]{OpenAIPacket2026}}]
\label{thm:packet}
For a fixed phase as in Equation~\eqref{eq:spherical-continuation},
every fixed admissible complexity range, $0<\kappa<1/10$, and
$\eta>0$, there is a finite constant $C$ such that
\begin{equation}\label{eq:packet-source}
 \mathcal A_m(T_md)^3
       \le C m^{10\kappa+\eta}G_{\delta,1}(d)^3
\end{equation}
for all its single-step maps and all its finite complex input
arrays.  The source also permits $A_{\delta,1}(d)^3$ on the right.
The constant may depend on $\kappa$, $\eta$, the complexity
range, the phase, and the windows.

In particular, suppose the initial labels
$\nu\in\delta\mathbb Z^2$ and
$z\in D_{\mathrm f}^{-1}\delta\mathbb Z^2$ lie in fixed bounded
sets.  Define
\begin{align}
 S d(\xi)
 &=(D_{\mathrm f}\delta)^{-2}
     \sum_{\nu,z}d_{\nu,z}\chi_\delta^\nu(\xi)
                               e(-N^2z\cdot\xi),
                     \label{eq:packet-synthesis}\\
 c_{\nu',z',j}
 &=\int Sd(\xi)\chi_1^{\nu'}(\xi)
                 e\bigl(N^2z'\cdot\xi+j\phi(\xi)\bigr)\,d\xi,
                     \label{eq:packet-samples}
\end{align}
where $\nu'\in\mathbb Z^2$,
$z'\in D_{\mathrm f}^{-1}N^{-2}\mathbb Z^2$, and
$j=0,\ldots,N^2-1$.  Retain any output coordinate mask for which
$\nu',z'$ lie in fixed bounded sets.  There is a finite
$P(\kappa,\eta)\ge1$, depending on these fixed sets, the phase,
and the windows, such that, for every dyadic $N\ge2$ and
$0<\eta\le1$,
\begin{equation}\label{eq:packet-bound}
 N^{-6}\sum_{\nu',z',j}|c_{\nu',z',j}|^3
 \le P(\kappa,\eta)N^{10\kappa+\eta}
                  m(d)\mathcal K_\kappa(d)^{1/2}.
\end{equation}
The sum is over the retained coordinates.
\end{theorem}

\begin{proof}[Specialization of the cited theorem]
Set $m=N$.  The scales in Equation~\eqref{eq:packet-scales} become
\[
\begin{array}{c|ccc}
 &\theta&r&w\\ \hline
 \text{input}&\delta&1&\delta\\
 \text{output}&1&N^{-2}&N^{-2}
\end{array}
\qquad t_I=\frac{j}{N^2},\qquad rw^2=N^{-6}\quad\text{at output}.
\]
Thus the companion's synthesis and
analysis operators are exactly the displayed operators, with
its window constant $L_{\mathrm f}$ equal to our $D_{\mathrm f}$.
The initial gauge is exactly the last expression in
Equation~\eqref{eq:packet-bound}, and
Equation~\eqref{eq:packet-atomic-lower} gives
\[
 N^{-6}\sum_{\nu',z',j}|c_{\nu',z',j}|^3
      \le\mathcal A_N(c)^3.
\]
All bounded label sets fit one fixed admissible complexity range
for sufficiently large $N$, since $m=N$; the finitely many
smaller dyadic scales may be included by enlarging the constant.
The bound at any such fixed scale follows from finite-dimensional
boundedness and $\mathcal K_\kappa(d)\ge\|d\|_\infty^2$.
Applying Equation~\eqref{eq:packet-source} proves
Equation~\eqref{eq:packet-bound}.  This argument uses arbitrary
input arrays, so it also applies after any restriction of their
support.
\end{proof}

Theorem~\ref{thm:packet} is the additional oscillatory input.  Its
source proves Equation~\eqref{eq:packet-source} in
\cite[Section~\PacketProofSection]{OpenAIPacket2026}; no rate for the dependence
of $P(\kappa,\eta)$ on its small parameters is asserted here.
